\chapter{Parity Philosophy and Permanent Deviations}
\label{ch:parity-philosophy}

Every project this book has covered draws a line between what it will faithfully reproduce
and what it will not. \texttt{sharp-runtime}'s own line is the most explicitly reasoned of
any in this ecosystem --- its \texttt{CLAUDE.md} states not just which .NET features are
permanently out of scope, but why, in enough detail that this chapter can quote the reasoning
directly rather than paraphrase it.

\section{Reflection: correctly a stub, not a gap}

\cnaclass{System.Type}, \cnaclass{Activator}, and enum reflection (\texttt{Enum.GetNames}/%
\texttt{GetValues}) are, in the project's own words, ``completely out of scope. Stubs are the
correct end state.'' This is a stronger claim than ``not yet implemented'' --- it states that
a stub \emph{is} the intended final behavior, not a placeholder awaiting more engineering.
Reflection depends on runtime type metadata a compiled C\texttt{++} binary simply does not
carry the way a .NET assembly does; reproducing it would mean building a real, if partial,
runtime type system from scratch, disproportionate to what CNA's own XNA-facing code actually
needs.

\section{Garbage collection: no-ops, by design}

\cnaclass{System.GC}'s methods are all no-ops. Object lifetime in this ecosystem is managed
by RAII and \texttt{shared\_ptr}, consistently, project-wide --- there is no collector to
force a collection of, or to query generation statistics from, so the methods exist only for
API-surface completeness, matching real XNA/FNA code that might call them defensively without
actually depending on their effect.

\section{Delegates: a three-tier model, corrected once in the project's own history}

This is worth covering in detail because the project's own documentation records a real
self-correction: an earlier description of the delegate model was found inaccurate during an
audit and rewritten. The corrected, current model has three distinct tiers, each solving a
different shape of problem. Most delegate-shaped C\# types --- \texttt{Action},
\texttt{Func}, most \texttt{*Callback} aliases --- are bare \texttt{using} aliases over
\texttt{std::function}: single-target, no multicast, no \texttt{BeginInvoke}/%
\texttt{EndInvoke}, sufficient for the overwhelming majority of call sites that only ever
needed one callback anyway. \cnaclass{System::Delegate} is a real multicast base
(\texttt{Combine}, \texttt{Remove}, \texttt{GetInvocationList}) for the smaller set of call
sites that genuinely need C\#'s multicast-delegate combining semantics.
\cnaclass{System::MulticastAction<Args...>} (Chapter~\ref{ch:sharp-runtime-overview}) is a
third, purpose-built type specifically for multicast event \emph{fields}, using
token-based removal since C\texttt{++} callables have no target-plus-method equality to
compare by the way C\# delegates do. All three tiers agree on one further point:
\texttt{DynamicInvoke} always throws \texttt{NotImplementedException}, since invoking a
delegate by a runtime-resolved argument array is itself a reflection-adjacent capability
already ruled out above.

\section{Serialization, P/Invoke, and cryptography: three different reasons for the same conclusion}

Serialization is ``ignored, not needed for game code'' --- a scope judgment rather than a
technical impossibility. P/Invoke and general native interop are out of scope for the same
underlying reason Chapter~\ref{ch:design-philosophy}'s type-alias table exists in the first
place: this project already \emph{is} native code, so a marshaling layer for calling into
native code from managed code has no problem left to solve here.

Cryptography and TLS are the one deviation with the most fully spelled-out reasoning of all,
and it deserves quoting closely: symmetric and asymmetric cryptography, X.509 certificates,
and \cnaclass{SslStream} are out of scope by an explicit, dated project-owner decision, on the
grounds that ``implementing this correctly needs either a large new external dependency
(OpenSSL/mbedTLS) or a hand-rolled, security-critical implementation, neither of which is
worth it for game code.'' This is a materially different kind of reasoning than the other
deviations in this chapter --- not ``this feature has no meaning here'' (reflection, GC) and
not ``this is out of scope for capacity reasons'' (serialization), but an explicit refusal to
carry the security liability of a hand-rolled cryptographic implementation, or the dependency
weight of a real one, for a class of application that does not need transport security built
into its own runtime layer. Non-cryptographic hashing (MD5, the SHA family, HMAC, PBKDF2) is
explicitly kept in scope and unaffected by this decision, on the reasoning that ``no key
material, no confidentiality guarantees to get wrong'' --- a hash function used for checksums
or content addressing carries none of the security risk a real encryption implementation
would.

\section{Two further, narrower accepted gaps}

Two smaller, explicitly accepted limitations round out this chapter. A fixed list of
POSIX-only subsystems (\cnaclass{Net::Sockets}, \cnaclass{IO::RandomAccess},
\texttt{AppDomain}/\texttt{AppContext}, \texttt{TimeZoneInfo}, \texttt{Diagnostics::Process},
POSIX signal handling, \texttt{NetworkInterface}, \texttt{FileSystemWatcher}) is accepted as
Linux/POSIX-only rather than cross-platform, for now. And \texttt{Decimal}, \texttt{Int128},
and \texttt{UInt128} are unsupported specifically on MSVC, since real 128-bit integer support
requires GCC or Clang's \texttt{\_\_int128} extension --- a deliberate, dated decision not to
hand-roll 128-bit arithmetic purely to cover one compiler, rather than an oversight.

\section{A renaming story worth knowing, as a model of how this project handles a breaking change}

One small, concrete episode illustrates how this project treats an intentional breaking API
change, and is worth contrasting directly with Chapter~\ref{ch:design-philosophy}'s ``no
backward-compatibility hacks'' rule in Volume~I: a \texttt{getCurrent()} accessor was renamed
to \texttt{getCurrentProperty()}, project-wide, across 23 files and 61 call sites, in one
single commit, with no backward-compatible alias left behind at all. The reasoning recorded
for this choice is direct: downstream consumers are expected to fix their own call sites, not
be shielded from a naming-convention correction by a compatibility shim that would itself
become permanent debt. This is the same principle Chapter~\ref{ch:design-philosophy}
described for CNA's own \cnaclass{Color::CornflowerBlue} example, applied here at a larger
scale.

\section{CI: an honest, offered, and declined gap}

One further fact is worth stating for the same reason this book has stated CNA's own CI scope
honestly (Chapter~\ref{ch:building-cna} in Volume~I): \texttt{sharp-runtime} has no tracked,
automated CI pipeline at all. This was itself a specific finding in the post-stabilization
audit from Chapter~\ref{ch:sharp-runtime-namespaces}, explicitly offered to the project owner
as something to address, and explicitly declined --- a local pre-push script
(\texttt{scripts/local\_ci\_check.sh}) serves as the equivalent gate instead, configuring,
building with a zero-warnings-and-errors requirement, and running the full test suite before
every push. Cross-compilation to Windows (MinGW-w64) and to Emscripten is claimed working, per
earlier stabilization work, but the test binary itself has never actually been cross-compiled
for either target in this project's own history --- GoogleTest is not cross-compiled for
either wasm or MinGW here --- so ``the library cross-compiles'' and ``the test suite passes
under cross-compilation'' remain two distinct, only-partially-verified claims, stated as such
rather than conflated.
